\documentclass{article}
\usepackage[utf8]{inputenc}
\usepackage{tabularx}
\usepackage{booktabs}
%\usepackage{times}
%\usepackage{newtxmath}
\PassOptionsToPackage{prologue,dvipsnames}{xcolor}


%\usepackage[mathscr]{euscript}
\usepackage{amssymb,setspace,lipsum,enumerate,amsthm,amsmath,mathtools,varwidth,calligra,hyperref,mathrsfs,stmaryrd,mathabx}
\usepackage{quiver}

\usepackage{titlesec}
\titleformat{\section}{\normalfont\scshape\filcenter}
  {\thesection.}
  {0.5 em}
  {}

\titleformat{\subsection}[runin] % run-in style, no line break
  {}                    % formatting of the title
  {\thesubsection.}               % label (number)
  {0.5em}                         % space between number and title
  {\bfseries}                               % before-code (empty)
  [.]   

\usepackage[margin=1.75in]{geometry}

\usepackage{tocloft}
\renewcommand{\cfttoctitlefont}{\normalfont\scshape}% Remove \bfseries from ToC title
\renewcommand{\cftsecfont}{}% Remove \bfseries from section titles in ToC
\renewcommand{\cftsecpagefont}{}% Remove \bfseries from section titles' page in ToC
\renewcommand{\contentsname}{\hfill\normalfont\scshape Contents\hfill}   
\renewcommand{\cftaftertoctitle}{\hfill}

% =============================================

\date{}
\author{\normalsize\scshape Emerson Hemley}
% Reduce space before/after title
%\title{\vspace{-1.5cm}\large\textbf{The Gauss-Manin connection via the de Rham space}\vspace{-0.5cm}}

\title{\vspace{-1.5cm}\large\textbf{The Gauss-Manin connection via the de Rham space}\vspace{-0.5cm}}


% =============================================

% Redefine abstract environment
\renewenvironment{abstract}{%
  \par\noindent{\textsc{Abstract.}}\small }
{\par\vspace{0.2em}}

% =============================================

\newtheorem{theorem}{Theorem}[section]
\newtheorem{lemma}[theorem]{Lemma}
\newtheorem{proposition}[theorem]{Proposition}
\newtheorem{corollary}[theorem]{Corollary}
\newtheorem{conjecture}[theorem]{Conjecture}
\theoremstyle{definition}
\newtheorem{definition}[theorem]{Definition}
\newtheorem{example}[theorem]{Example}
\newtheorem{question}[theorem]{Question}
\newtheorem{exercise}[theorem]{Exercise}
\theoremstyle{remark}
\newtheorem{remark}[theorem]{Remark}
\newtheorem{caution}[theorem]{Caution}

% =============================================

%\DeclareMathOperator{\HOM}{\mathscr{H}  \mathscr{O} \mathscr{M}}
\DeclareMathOperator{\HOM}{\mathscr{H}\kern -3pt \mathscr{O} -3pt \mathscr{M}}
\newcommand{\scr}[1]{\mathscr{#1}}
\newcommand{\frk}[1]{\mathfrak{#1}}
\newcommand{\mc}[1]{\mathcal{#1}}
\newcommand{\id}{\textrm{id}}
\newcommand{\maps}{\text{Maps}}
\DeclareMathOperator{\spec}{Spec}
\DeclareMathOperator{\spp}{sp}
\DeclareMathOperator{\sym}{Sym}
\DeclareMathOperator{\homm}{Hom}
\renewcommand{\hom}{\homm}
\renewcommand{\arraystretch}{1.5}
\DeclareMathOperator{\qcoh}{QCoh}
\DeclareMathOperator{\coh}{Coh}
\DeclareMathOperator{\perf}{Perf}
\DeclareMathOperator{\vect}{Vect}
\DeclareMathOperator{\loc}{Loc}
\renewcommand{\sslash}{\mathbin{/\mkern-6mu/}}
\DeclareMathOperator{\z}{\mathbb{Z}}
\DeclareMathOperator{\q}{\mathbb{Q}}
\DeclareMathOperator{\im}{im}
\DeclareMathOperator{\gal}{Gal}
\DeclareMathOperator{\End}{End}
\DeclareMathOperator{\br}{Br}
\DeclareMathOperator{\pic}{Pic}
\DeclareMathOperator{\dr}{dR}
\DeclareMathOperator{\dol}{Dol}
\DeclareMathOperator{\crys}{crys}
\DeclareMathOperator{\an}{an}
\newcommand{\ds}{\displaystyle}
\DeclareMathOperator{\pr}{pr}
\DeclareMathOperator{\red}{red}
\DeclareMathOperator{\corr}{Corr}
\DeclareMathOperator{\cat}{Cat}
\DeclareMathOperator{\cech}{\check{C}ech}
\DeclareMathOperator{\h}{H}
\DeclareMathOperator{\dqc}{D_{qc}}
\DeclareMathOperator{\aff}{Aff}
\DeclareMathOperator{\algstk}{AlgStk}
\DeclareMathOperator{\op}{op}
\DeclareMathOperator{\set}{Set}

\newcommand{\boxx}[1]{\begin{center}
\fbox{\begin{varwidth}{\dimexpr\textwidth-2\fboxsep-2\fboxrule\relax}{#1}
\end{varwidth}}
\end{center}}
\usepackage[OT2,T1]{fontenc}
\DeclareSymbolFont{cyrletters}{OT2}{wncyr}{m}{n}
\DeclareMathSymbol{\sh}{\mathalpha}{cyrletters}{"58}

\usepackage[dvipsnames]{xcolor}

\renewcommand{\labelenumi}{(\alph{enumi})} 

\usepackage{relsize}
\usepackage[bbgreekl]{mathbbol}
\usepackage{amsfonts}
\DeclareSymbolFontAlphabet{\mathbb}{AMSb} %to ensure that the meaning of \mathbb does not change
\DeclareSymbolFontAlphabet{\mathbbl}{bbold}
\newcommand{\Prism}{{\mathlarger{\mathbbl{\Delta}}}}

\begin{document}

\maketitle

\begin{abstract}
Let $k$ be a field of characteristic 0 and $X \to S$ be a smooth morphism of smooth $k$-schemes. It is a classical result that the relative algebraic de Rham cohomology $\h_{\dr}^i(X/S)$ is naturally endowed with a flat connection, the \emph{Gauss-Manin} connection. In this note, we provide an alternative perspective coming from the de Rham space. In particular, we explain how the relative de Rham cohomology (considered as a $\scr{D}_S$-module under the Gauss-Manin connection) is isomorphic to the higher direct image of $\scr{O}_{X_{\dr}}$ along the induced  morphism between de Rham spaces $X_{\dr} \to S_{\dr}$.

\end{abstract}

\tableofcontents

\section{The de Rham space}

\subsection{Preliminaries} We introduce the de Rham space, originally defined in \cite{Simpson1996}. Many of the technical arguments in this section are found in \cite{Rib24}. 

Let $k$ be a fixed field of characteristic 0 and $\text{Aff}_k$ be the category of affine $k$-schemes. Let $\tau$ be a Grothendieck topology on $\text{Aff}_k$. We define a $\tau$ algebraic stack to be a sheaf $\text{Aff}_k \to \text{Grpd}$ in the $\tau$ topology. The category of $\tau$ algebraic stacks will be denoted $\algstk_\tau$.

\begin{definition}[de Rham space] Given any presheaf of $X: \text{Aff}_k \to \set$, we define $X_{\dr}$ to be the presheaf 
$$ X_{\dr}(A) \coloneqq X(A^{\red})$$
Note that there is a natural map $X \to X_{\dr}$ and that associated to a morphism $X \to S$ there is an induced morphism $X_{\dr} \to S_{\dr}$. In this case, we define the relative de Rham space $(X/S)_{\dr}$ to be the fibered product $X_{\dr} \times_{S_{\dr}} S$.
\end{definition}


Let $X$ be a scheme which is formally smooth over $k$. If a $k$-algebra $A$ has nilpotent nilradical (for instance, if $A$ is Noetherian), then every $A^{\red}$-point of $X$ extends to an $A$-point,
\[\begin{tikzcd}
  {\spec A^{\red}} & X \\
  {\spec A} & {\spec k}
  \arrow[from=1-1, to=1-2]
  \arrow[from=1-1, to=2-1]
  \arrow[from=1-2, to=2-2]
  \arrow[dashed, from=2-1, to=1-2]
  \arrow[from=2-1, to=2-2]
\end{tikzcd}\]
by infinitesimal lifting.
Consequently, the natural map $X \to X_{\dr}$ is an epimorphism of presheaves. In general though, $X$ smooth over $k$ is sufficient. 

\begin{proposition} If $X$ is smooth over $k$, then $X \to X_{\dr}$ is an epimorphism.
\end{proposition} 
\begin{proof}

Let $I_{n}$ be the ideal of $A$ generated by elements of the form $x^n = 0$, and define $A_n \coloneqq A/I_n$. Then $ A^{\red} \cong \varinjlim_n A_n$, and
$$ \hom(\spec A^{\red}, X) \cong \hom(\varprojlim_n \spec A_n, X) \cong \varinjlim_n \hom(\spec A_n,X)$$
where the second isomorphism holds because because $X$ is locally finite presentation \cite{stacks-project} \href{https://stacks.math.columbia.edu/tag/01ZC}{Tag 01ZC}. Consider an element on the right; this is given by a representative $\spec A_n \to X$ for some $n$. By formal smoothness, this lifts to $\spec A \to X$.
\end{proof}

In view of the above proof, we also have the following characterization: if $X$ is locally of finite presentation over $k$, then
$$ X_{\dr}(A) \cong \hom(\varprojlim_n \spec A_n,X) \cong \varinjlim_n \hom(\spec A_n,X)$$
and therefore $X_{\dr}$ is a colimit. An important consequence of this is the following.

\begin{corollary} If $X$ is locally of finite presentation over $k$, $X_{\dr}: \aff_k \to \set$ commutes with finite limits and arbitrary colimits.
\end{corollary}


An important property of the de Rham space is that it satisfies descent.

\begin{proposition}[\cite{Rib24}, proposition 1.1.14] If $X$ is a locally finite type scheme, the de Rham space $X_{\dr}$ satisfies fppf descent.
\end{proposition}

We emphasize that the assumption $k$ is of characteristic 0 is necessary. Thus by our conventions, $X_{\dr}$ is an fppf algebraic stack. 

\subsection{Algebraic de Rham cohomology} A key feature of the de Rham space is that the cover $X \to X_{\dr}$ provides an identification of the algebraic de Rham cohomology of $X$ with the $\scr{O}$-coherent cohomology of $X_{\dr}$. We will provide a proof a relative version of this theorem. 

Let $f: X \to S$ be a morphism of $k$-schemes, and consider the induced morphisms
% https://q.uiver.app/#q=WzAsNCxbMCwxLCJTIl0sWzEsMCwiWF97XFxkcn0iXSxbMSwxLCJTX3tcXGRyfSJdLFswLDAsIihYL1MpIl0sWzEsMiwiZl97XFxkcn0iXSxbMCwyLCJcXHBpIiwyXSxbMywwLCJcXHdpZGV0aWxkZXtmfSIsMl0sWzMsMSwiXFx3aWRldGlsZGV7XFxwaX0iXV0=
\[\begin{tikzcd}
  {(X/S)_{\dr}} & {X_{\dr}} \\
  S & {S_{\dr}}
  \arrow[, from=1-1, to=1-2]
  \arrow["{\widetilde{f}_{\dr}}"', from=1-1, to=2-1]
  \arrow["{f_{\dr}}", from=1-2, to=2-2]
  \arrow[, from=2-1, to=2-2]
\end{tikzcd}\]
Let $\scr{O}$ be the structure sheaf of $(X/S)_{\dr}$.

\begin{theorem}\label{dR-comp} If $f: X \to S$ is a morphism of smooth $k$-schemes,
$$ \h^i_{\dr}(X/S) \cong R^i \widetilde{f}_{\dr*} \scr{O}$$
where the left is the Zariski hypercohomology of $\Omega^\bullet_{X/S}$.
\end{theorem}

In order to prove \ref{dR-comp}, we need the following lemmas.

\begin{lemma}\label{formal-completion} Let $X$ be locally of finite presentation. Then, the formal completion of $X$ along a closed subscheme $Z$ may be identified as
$$ \widehat{X}_{Z} \cong X \times_{X_{\dr}} Z_{\dr} $$
\end{lemma}

\begin{lemma}\label{diagonal} If $X$ and $S$ are locally of finite presentation, 
$$ X \times_{(X/S)_{\dr}} X \cong \widehat{(X \times_S X)}_\Delta$$
where the right is the formal completion of the image of the diagonal $X \to X \times_S X$.
\end{lemma}

Some remarks are in order. A priori, the definition of formal completion of $X$ along a subscheme $Z$ requires $Z$ to be closed. However in \ref{formal-completion}, note that the left hand side does not require $Z \to X$ to be a closed immersion. Thus in \ref{diagonal}, we do not require $X \to S$ to be separated and instead define formal completion by the above formula.

\begin{proof}

Let $Y = X \times_S X$, so that $\widehat{Y}_\Delta \cong Y \times_{Y_{\dr}} X_{\dr}$. Then the fact follows from% https://q.uiver.app/#q=WzAsNCxbMSwwLCJYIFxcdGltZXNfUyBYIl0sWzAsMSwiWF97XFxkcn0iXSxbMSwxLCIoWFxcdGltZXNfUyBYKV97XFxkcn0iXSxbMCwwLCJYX3tcXGRyfSBcXHRpbWVzX3tTX3tcXGRyfX0gUy{}JdLFsxLDIsIlxcRGVsdGEiLDJdLFswLDJdLFszLDFdLFszLDBdXQ==
\[\begin{tikzcd}
  {X_{\dr} \times_{S_{\dr}} S} & {Y = X \times_S X} \\
  {X_{\dr}} & {Y_{\dr} \cong X_{\dr} \times_{S_{\dr}} X_{\dr}}
  \arrow[from=1-1, to=1-2]
  \arrow[from=1-1, to=2-1]
  \arrow[from=1-2, to=2-2]
  \arrow["\Delta"', from=2-1, to=2-2]
\end{tikzcd}\]
being cartesian.
\end{proof}

\begin{proof}[Proof of \ref{dR-comp}] 

For simplicity, we consider the absolute case $S = \spec k$ first. Since $\widehat{X \times X} \rightrightarrows X \to X_{\dr}$ is a coequalizer of presheaves, $X \to X_{\dr}$ is an effective epimorphism. Therefore the simplicial scheme $\widehat{X}^\bullet = \cech(X \to X_{\dr})$ is a hypercover of $X_{\dr}$. By \ref{diagonal}, $\widehat{X}^n$ is the ind-scheme given by formal completion of the $n$-fold product $X \times \cdots \times X$ along the diagonal. By \cite{Rod25}, there is a quasi-isomorphism
$$ \varprojlim \scr{O}_{\widehat{X}^\bullet} \to \Omega^\bullet_{X/k}$$
of Zariski sheaves on $X$, and consequently
$$ R\Gamma(X_{\dr},\scr{O}_{X_{\dr}}) \cong R\Gamma(\widehat{X}^\bullet, \scr{O}_{\widehat{X}^\bullet}) \cong R\Gamma(X,\Omega^\bullet_{X/k})$$
finishes the proof. The relative case $X \to S$ follows the same way.
\end{proof}
{}
%\subsection{Positive characteristic}

\subsection{Quasi-coherent sheaves on de Rham spaces} 
A key insight of Simpson was that for a smooth $k$-schemes $X$, that the category of quasi-coherent sheaves over $X_{\dr}$ identifies with the category of algebraic $\scr{D}$-modules on $X$. By \ref{diagonal} and \ref{formal-completion},
$$ \widehat{X \times X} \rightrightarrows X \to X_{\dr} $$
is a coequalizer of presheaves and thus we can view $X_{\dr}$ as a quotient of $X$ by the relation of \emph{infinitesimal closeness} (c.f \cite{Lur}). The data of a quasi-coherent sheaf on $X_{\dr}$ is thus equivalent to a quasi-coherent sheaf on $X$ along with a compatible system of isomorphism over points which lie in some $n$-fold formal thickening of the diagonal. This is equivalent to an action of the ring of differential operators on $X$.

\begin{theorem}\label{equiv1} If $X$ is a smooth $k$-scheme, the category of quasi-coherent sheaves on $X_{\dr}$ is equivalent to the category of quasi-coherent $\scr{D}_X$-modules. That is,
$$\qcoh(X_{\dr}) \cong \qcoh(\scr{D}_X)$$
\end{theorem}

See also \cite{Rod25} \S 9. In particular, a locally free sheaf on $X_{\dr}$ is equivalent to a locally free sheaf along with a flat connection on $X$. Pulling back a locally free sheaf along $X \to X_{\dr}$ corresponds to forgetting the flat connection. Therefore, the fact that the relative de Rham cohomology sheaf $\h^i_{\dr}(X/S)$ may be endowed with a flat connection must correspond to the fact that it is pulled back from some quasi-coherent sheaf on $S_{\dr}$.


Consider the following diagram,
% https://q.uiver.app/#q=WzAsNCxbMSwwLCJYX3tkUn0iXSxbMSwxLCJTX3tkUn0iXSxbMCwxLCJTIl0sWzAsMCwiKFgvUylfe2RSfSJdLFszLDAsImcnIl0sWzMsMiwiZidfe2RSfSIsMl0sWzIsMSwiZyIsMl0sWzAsMSwiZl97ZFJ9Il1d
\[\begin{tikzcd}
  {(X/S)_{\dr}} & {X_{\dr}} \\
  S & {S_{\dr}}
  \arrow["{\widetilde{\pi}}", from=1-1, to=1-2]
  \arrow["{\widetilde{f}_{\dr}}"', from=1-1, to=2-1]
  \arrow["\pi"', from=2-1, to=2-2]
  \arrow["{f_{\dr}}", from=1-2, to=2-2]
\end{tikzcd}\tag{$\ast$}\]
\ref{dR-comp} asserts that the relative de Rham cohomology is isomorphic to $R^i \widetilde{f}_{\dr*} \scr{O}$. Recall that there is a natural map
\[\pi^* R^i f_{\dr*} \scr{O}_{X_{\dr}} \to R^i \widetilde{f}_{\dr*} \widetilde{\pi}^* \scr{O}_{X_{\dr}}
\] 
called the base-change morphism. If this map is an isomorphism, observe that this would formally endow the relative de Rham cohomology with a flat connection: since $R^i f_{\dr*} \scr{O}_{X_{\dr}}$ is a quasi-coherent sheaf on $S_{\dr}$, it is equipped with one by \ref{equiv1}. The classical construction of the Gauss-Manin connection occurs because we were going ``the wrong way'' around the above cartesian square.


\begin{proposition} Let $f: X \to S$ be a smooth morphism of smooth $k$-schemes. If $(\ast)$ satisfies base-change, there exists a canonical flat connection $\nabla$ on $\h^i_{\dr}(X/S)$.
\end{proposition}


\section{The six functor formalism for $X_{\dr}$}

 Because of the non-algebraic nature of $S_{\dr}$, establishing the above base change is slightly delicate; since $S_{\dr}$ is not a scheme, we cannot apply traditional proper or smooth base change for quasi-coherent sheaves on schemes. To this end, we proceed by setting up some abstract machinery: we appeal to the six functor formalism for quasi-coherent sheaves on $X_{\dr}$ as found in \cite{Rib24} or \cite{Rod25}.



\subsection{Abstract 6-functor formalisms}

We recall the definition of a 6-functor formalism. Rather than strive for an exhaustive treatment, we will introduce only the minimal amount of theory to state the desired result. For details see \cite{Sch22}.


\begin{definition} A \emph{geometric set up} a pair $(\mathcal{C}, E)$ where $\mathcal{C}$ is an $\infty$-category with finite limits and $E$ is a collection of morphisms in $\mathcal{C}$ containing all isomorphisms and stable under pull back, composition, and taking diagonals.
\end{definition}

\begin{remark}
$E$ is the class of $!$-able morphisms, or morphisms which admit a $f_!$.
\end{remark}

\begin{definition} Given a geometric set up $(\mathcal{C},E)$, the category of correspondences $\corr(\mathcal{C},E)$ is defined to be the category whose objects are the objects of $\mathcal{C}$ and a morphism from $X$ to $Y$ is a diagram
% https://q.uiver.app/#q=WzAsMyxbMSwwLCJaIl0sWzAsMSwiWCJdLFsyLDEsIlkiXSxbMCwyXSxbMCwxXV0=
\[\begin{tikzcd}
  & Z \\
  X && Y
  \arrow[from=1-2, to=2-1]
  \arrow[from=1-2, to=2-3]
\end{tikzcd}\]
in $\mathcal{C}$ where the arrow on the right $Z \to Y$ lies in $E$. A composition of morphisms is
% https://q.uiver.app/#q=WzAsNixbMCwyLCJYIl0sWzEsMSwiWiJdLFsyLDIsIlkiXSxbMywxLCJTIl0sWzQsMiwiVCJdLFsyLDAsIlogXFx0aW1lc19ZUyJdLFsxLDBdLFsxLDJdLFszLDJdLFszLDRdLFs1LDNdLFs1LDFdXQ==
\[\begin{tikzcd}
  && {Z \times_YS} \\
  & Z && S \\
  X && Y && T
  \arrow[from=1-3, to=2-2]
  \arrow[from=1-3, to=2-4]
  \arrow[from=2-2, to=3-1]
  \arrow[from=2-2, to=3-3]
  \arrow[from=2-4, to=3-3]
  \arrow[from=2-4, to=3-5]
\end{tikzcd}\]
where the left arrow $Z \times_Y S \to T$ lies in $E$ by assumption. There is a monoidal structure on $\corr(\mathcal{C},E)$ given by the product in $\mathcal{C}$.
\end{definition}


\begin{definition} Given a geometric set up $(\mathcal{C},E)$, a 3-functor formalism is a lax symmetric monoidal functor $\mathcal{D}: \corr(\mathcal{C},E) \to \cat_\infty$.
\end{definition}

Given an arbitrary morphism $f: X \to Y$ in $\mathcal{C}$, the image of the correspondence
% https://q.uiver.app/#q=WzAsMyxbMSwwLCJYIl0sWzAsMSwiWSJdLFsyLDEsIlgiXSxbMCwyLCJcXHRleHR7aWR9Il0sWzAsMSwiZiIsMl1d
\[\begin{tikzcd}
  & X \\
  Y && X
  \arrow["f"', from=1-2, to=2-1]
  \arrow["{\text{id}}", from=1-2, to=2-3]
\end{tikzcd}\]
is defined to be $f^*: \mathcal{D}(Y) \to \mathcal{D}(X)$. If $f: X \to Y$ lies in $E$, the image of the correspondence
% https://q.uiver.app/#q=WzAsMyxbMSwwLCJYIl0sWzAsMSwiWCJdLFsyLDEsIlkiXSxbMCwyLCJmIl0sWzAsMSwiXFx0ZXh0e2lkfSIsMl1d
\[\begin{tikzcd}
  & X \\
  X && Y
  \arrow["{\text{id}}"', from=1-2, to=2-1]
  \arrow["f", from=1-2, to=2-3]
\end{tikzcd}\]
is defined to be $f_!: \mathcal{D}(X) \to \mathcal{D}(Y)$.

\begin{definition} A 6-functor formalism is a 3-functor formalism
$$ \mathcal{D}: \corr(\mathcal{C},E) \to \cat_\infty$$
such that the following holds,
\begin{enumerate}
  \item the functors $f^*$ and $f_!$ both admit right adjoints (defined to be $f_*$ and $f^!$),
  \item for each object of $\mathcal{C}$, the symmetric monoidal $\infty$-category $\mathcal{D}(X)$ is closed.
\end{enumerate}
\end{definition}

A formal consequence of this is the following: given any Cartesian diagram
% https://q.uiver.app/#q=WzAsNCxbMCwwLCJYJyJdLFsxLDAsIlgiXSxbMSwxLCJZIl0sWzAsMSwiWSciXSxbMCwzLCJmJyIsMl0sWzEsMiwiZiJdLFszLDIsImciLDJdLFswLDEsImcnIl1d
\[\begin{tikzcd}
  {X'} & X \\
  {Y'} & Y
  \arrow["{g'}", from=1-1, to=1-2]{}
  \arrow["{f'}"', from=1-1, to=2-1]
  \arrow["f", from=1-2, to=2-2]
  \arrow["g"', from=2-1, to=2-2]
\end{tikzcd}\]
in $\mathcal{C}$ where $f$ is $!$-able, 
there is a specified isomorphism of functors $$g^* \circ f_! \to f'_! \circ (g')^*$$ \cite{Mann22}. This is a salient feature of 6-functor formalisms and the motivation for our consideration of them.

\begin{definition} Let $f: X \to Y$ be $!$-able and $\omega_f \coloneqq f^! (1_Y)$ be the dualizing object. The morphism $f$ is said to be \emph{cohomologically smooth} provided that
\begin{enumerate}
  \item The natural transformation $\omega_f \otimes f^*(-) \to f^!(-)$ is an equivalence,
  \item  $\omega_f$ is $\otimes$-invertible in $\mathcal{D}(X)$,
  \item For any morphism $g: Y' \to Y$ with base change $f': X' \to Y'$ of $f$, the above properties also hold for $f'$. Moreoever, if $g': Y' \to Y$ denotes the base change of $g$, we also require the natural map $$(g')^* \omega_f \to \omega_{f'}$$ to be an isomorphism.
\end{enumerate}
\end{definition}
For more details see \cite{Sch22} \S 5. Cohomological smoothness is important because it yields abstract notion of smooth base change (\cite{Zavyalov2023}, lemma 2.3.9).
\begin{theorem}[Cohomologically smooth base change]\label{smooth base change} Let
% https://q.uiver.app/#q=WzAsNCxbMCwxLCJZJyJdLFsxLDAsIlgiXSxbMSwxLCJZIl0sWzAsMCwiWCciXSxbMSwyLCJmIl0sWzAsMiwiZyIsMl0sWzMsMCwiZiciLDJdLFszLDEsImcnIl1d
\[\begin{tikzcd}
  {X'} & X \\
  {Y'} & Y
  \arrow["{g'}", from=1-1, to=1-2]
  \arrow["{f'}"', from=1-1, to=2-1]
  \arrow["f", from=1-2, to=2-2]
  \arrow["g"', from=2-1, to=2-2]
\end{tikzcd}\]
be a Cartesian square in $\mathcal{C}$. Then the natural morphism
$$g^* \circ f_* \to (f')_* \circ (g')^*$$
is an isomorphism if $g$ is cohomologically smooth.
\end{theorem}

As a final point, we mention briefly a particular type of cohomologically smooth morphism called cohomological étale morphisms. Intuitively, cohomological étale morphisms are cohomologically smooth with trivial dualizing object, i.e. satisfying $f^* \cong f^!$. We do not include a precise definition since cohomological smoothness is a sufficient condition for our needs  (c.f. \cite{Zavyalov2023} for a precise definition).


\subsection{Quasi-coherent sheaves on algebraic stacks}

We outline a construction for a 6-functor formalism for the derived category of quasi-coherent sheaves on algebraic stacks. This is relies on subtle extension theorems for 6-functor formalims which we will treat as a blackbox. This is done in full detail in \cite{Rod25} and \cite{Rib24}.

\smallskip

As before, let $\aff_k$ be the category of affine $k$-schemes.

\begin{theorem} The functor $\dqc: \aff_k^{\op} \to \cat_{\infty}$ which maps $X$ to its derived $\infty$-category of quasi-coherent sheaves $\dqc(X)$ extends to a 6-functor formalism on the geometric setup $(\aff_k, \emph{all})$ consisting of all morphisms. 
\end{theorem}

\begin{proof} This is proved in \cite{Sch22} \S 8.3. See also \cite{Rib24} theorem 1.3.20 or \cite{Rod25} proposition 6.11.
\end{proof}

To extend this to a 6 functor formalism on algebraic stacks, we appeal to general extension results of \cite{HM24}. This uses the $\dqc$-topology, though we will not elaborate on this since we can restrict to the coarser fppf topology.

\begin{theorem} The six functor formalism on the geometric set up $(\aff_k, \emph{all})$ extends uniquely to a six functor formalism $(\algstk_{\dqc}, E)$ such that
\begin{enumerate}
  \item Let $f : Y \to X$ be a map whose pullback to every object in $\aff_k$ lies in $E$, then $f$ lies in $E$.
  \item Let $f: Y \to X$ be a map which is $!$-locally on the source or target in $E$; then f lies in $E$.
  \item Every map $f: Y \to X$ in $E$ with $X \in $ Sch is $!$-locally on the source in $E$.
\end{enumerate}
\end{theorem}
\begin{proof}
\cite{Rod25}, Corollary 6.19.
\end{proof}

We emphasize that we do not have explicit knowledge of the class $E$. To check that a particular morphism of algebraic stacks is $!$-able, we need to verify (a)-(c). For an expository account, see \cite{Sch22} \emph{Appendix to Lecture IV: Passage to stacks}. 

\begin{remark}
The above result uses the $\dqc$-topology on algebraic stacks. However this does not pose a substantial problem; the $\dqc$-topology is \emph{finer} than the fppf topology, and \cite{Rib24} shows that there is a unique six functor formalism on fppf algebraic stacks whose $!$-able morphisms are precisely the morphisms whose sheafifications lie in the $E$ above (proposition 1.3.23). That is, there is a compatible 6-functor formalism on the geometric set up $(\algstk_{\text{fppf}}, E')$. In the sequel, we restrict to this formalism.
\end{remark}

We need the following technical lemmas to proceed. Let $X$ and $Y$ be schemes.

\begin{lemma}[\cite{Sch22}, Proposition 8.14]\label{proper=smooth} If $f: X \to Y$ is smooth and proper, then $f$ is cohomologically smooth.
\end{lemma}

\begin{lemma}[\cite{Rib24}, Proposition 1.3.25]\label{closed=smooth} Let $Z \subset X$ be a closed subscheme, where the ideal sheaf is locally finitely generated. Then the morphism
$$ \widehat{X}_Z \to X$$
is cohomologically étale.
\end{lemma}

\begin{proposition}\label{coh-smooth} If $X$ is locally finite type, $X \to X_{\dr}$ is cohomologically smooth.
\end{proposition}

We remark that the proof of this result is essentially (\cite{Sch22}, proposition 8.23), but Scholze works in the IndCoh-formalism.

\begin{proof} Cohomological smoothness can be verified locally on the source and target. Since this is local on the source, we can restrict to an affine $X$. By taking a compactification, we can restrict to a proper $X$. Let $Y \to X_{\dr}$ be an arbitrary map from a scheme. Then the following diagram is Cartesian, 
% https://q.uiver.app/#q=WzAsNCxbMCwxLCJYIl0sWzEsMSwiWF97XFxkcn0iXSxbMSwwLCJZIl0sWzAsMCwiKFggXFx0aW1lc19rIFkpX3tcXHdpZGVoYXR7Wn19Il0sWzAsMV0sWzIsMV0sWzMsMF0sWzMsMl1d
\[\begin{tikzcd}
  {\widehat{(X \times_k Y)}_{Z}} & Y \\
  X & {X_{\dr}}
  \arrow[from=1-1, to=1-2]
  \arrow[from=1-1, to=2-1]
  \arrow[from=1-2, to=2-2]
  \arrow[from=2-1, to=2-2]
\end{tikzcd}\]
where $Z \subset X \times_k Y$ is the graph of $Y_{\text{red}} \to X$, where $Z$ is closed because $X$ is separated and $Y^{\red} \to Y$ is a closed immersion. Then $\widehat{(X \times_k Y)}_{Z} \to Y$ is the composite of $\widehat{(X \times_k Y)}_{Z} \to X \times_k Y$ and $X \times_k Y \to Y$. Note that the first is cohomologically smooth by \ref{closed=smooth} and that the second is a base change of $X \to \spec k$ and is thus cohomologically smooth by \ref{proper=smooth}.
\end{proof}

\begin{proposition}
Let $f: X \to Y$ be a morphism between finite type schemes. Then $f_{\dr}: X_{\dr} \to Y_{\dr}$ is $!$-able in the 6-functor formalism for quasi-coherent sheaves.
\end{proposition}

\begin{proof}
Since having $!$-functors is local on the source, if suffices to verify this after precomposing with the cohomologically smooth map $X \to X_{\dr}$. Then $X \to Y_{\dr}$ is the composite of $f$ and $Y \to Y_{\dr}$, which are both $!$-able.
\end{proof}

\subsection{The Gauss-Manin connection}

Let $f: X \to S$ be a morphism of smooth $k$-schemes and assume that the induced map $f_{\dr}: X_{\dr} \to S_{\dr}$ admits $!$-functors (e.g. if $X$ and $S$ are finite type over $k$). Recall the diagram
% https://q.uiver.app/#q=WzAsNCxbMSwwLCJYX3tkUn0iXSxbMSwxLCJTX3tkUn0iXSxbMCwxLCJTIl0sWzAsMCwiKFgvUylfe2RSfSJdLFszLDAsImcnIl0sWzMsMiwiZidfe2RSfSIsMl0sWzIsMSwiZyIsMl0sWzAsMSwiZl97ZFJ9Il1d
\[\begin{tikzcd}
  {(X/S)_{\dr}} & {X_{\dr}} \\
  S & {S_{\dr}}
  \arrow["{\widetilde{\pi}}", from=1-1, to=1-2]
  \arrow["{\widetilde{f}}"', from=1-1, to=2-1]
  \arrow["\pi"', from=2-1, to=2-2]
  \arrow["{f_{\dr}}", from=1-2, to=2-2]
\end{tikzcd}\tag{$\ast$}\]
\ref{smooth base change} and \ref{coh-smooth} implies the following:

\begin{corollary} One has base-change for diagram $(*)$. In particular,
$$ \pi^*R^i f_{\dr *} \scr{O}_{X_{\dr}} \to R^i\widetilde{f}_* \widetilde{\pi}^* \scr{O}_{X_{\dr}}$$
is an isomorphism.
\end{corollary}


We emphasize the smoothness of $f$ is unnecessary to establish base change for the relative de Rham square; since $\pi$ is cohomologically smooth we need not assume any conditions on $f_{\dr}$ (other than that it's $!$-able). However, if $f$ is not smooth, then $\h^i_{\dr}(X/S)$ will not be a locally free sheaf and thus we speak of $\scr{D}_S$-module structures instead of connections.

\begin{corollary} There is a canonical $\scr{D}_S$-module structure on the relative de Rham cohomology $\h_{\dr}^i(X/S)$.
\end{corollary}


Moreover this $\scr{D}_S$-module is naturally isomorphic to the $i$-th higher direct image of the structure sheaf of $X_{\dr}$ along $f_{\dr}$.

%\newpage

\bigskip

\nocite{*}


% At the end of your document, before \end{document}
\bibliographystyle{alpha}  % or alpha, abbrv, etc.
{\small \bibliography{references}}

%{\footnotesize
%\bibliography{bibfile}}

\end{document}


\section{Appendix}

\subsection{The Kodaira-Spencer map}

\begin{proposition} The Higgs field on the relative Dolbeault cohomology $\h_{\dol}^i(X/S)$ given by cupping with the Kodaira-Spencer class is isomorphic to $R^i f_{\dol} \scr{O}_{\dol}$.
\end{proposition}


\subsection{Positive Characteristic}
Using Bhatt's definition of $X_{\dr}$ over $k = \mathbb{F}_p$, we expect the following: let $f:X \to \spec R$ be a morphism of smooth $k$-schemes.
\begin{proposition}
There is a canonical flat connection with nilpotent $p$-curvative on $H^i_{\crys}(X/W(R))$ isomorphic to $R^i f_{\dr *} \scr{O}_{X_{\dr}}$.
\end{proposition}
